An agent continuing a task after native context compaction may retain a
conclusion while losing the evidence needed to determine whether that
conclusion still applies. Reopening raw history can recover missing
information, but a relevant passage may describe an earlier code
revision, an incomplete test run, or an authorization that was later
withdrawn. This working draft studies whether explicit verification
improves subsequent decisions beyond strong raw-history retrieval. We
specify a selective verification layer that discovers candidate claims
from observable task state, retrieves and opens original sources without
evaluator-provided locations, and checks type-specific requirements for
test status, current versions, and historical authorization. Reuse of
previously verified evidence requires a certificate bound to a complete
archive view, a current archive watermark, the relevant object state,
and a valid time interval; a recent timestamp alone is insufficient. The
proposed evaluation pairs native continuations under a shared model,
archive, and evidence budget, comparing the method with Raw-RAG,
always-on verification, and reproducible neighboring approaches. Its
primary endpoint is task-level evidence-supported decision success. Cost
claims require both quality noninferiority and lower total standardized
processing cost. The draft makes no claim of empirical superiority: it
provides the method specification, falsifiable hypotheses, and evidence
required to determine when source verification is useful.
